\documentclass{article}

% if you need to pass options to natbib, use, e.g.:
%     \PassOptionsToPackage{numbers, compress}{natbib}
% before loading neurips_2019

% ready for submission
% \usepackage{neurips_2019}

% to compile a preprint version, e.g., for submission to arXiv, add add the
% [preprint] option:
%     \usepackage[preprint]{neurips_2019}

% to compile a camera-ready version, add the [final] option, e.g.:
     \usepackage[final]{style}

% to avoid loading the natbib package, add option nonatbib:
%     \usepackage[nonatbib]{neurips_2019}

\usepackage[utf8]{inputenc} % allow utf-8 input
\usepackage[T1]{fontenc}    % use 8-bit T1 fonts
\usepackage{hyperref}       % hyperlinks
\usepackage{url}            % simple URL typesetting
\usepackage{booktabs}       % professional-quality tables
\usepackage{amsfonts}       % blackboard math symbols
\usepackage{nicefrac}       % compact symbols for 1/2, etc.
\usepackage{microtype}      % microtypography

\title{Your Project Title}

% The \author macro works with any number of authors. There are two commands
% used to separate the names and addresses of multiple authors: \And and \AND.
%
% Using \And between authors leaves it to LaTeX to determine where to break the
% lines. Using \AND forces a line break at that point. So, if LaTeX puts 3 of 4
% authors names on the first line, and the last on the second line, try using
% \AND instead of \And before the third author name.

\author{%
  name1 \\
  Department \\
  \texttt{andrew email} \\
  \And
  name2 \\
  Department \\
  \texttt{andrew email} \\
  \And
  name3 \\
  Department \\
  \texttt{andrew email} \\
}

\begin{document}

\maketitle

\section{Problem Statement}
This is your project proposal! It is due on Friday, September 11th, 2026. You may use your late days on this assignment, but each group member will have to use one of their late days.

Be sure to check \url{https://piazza.com/cmu/fall2026/18794a/resources} for the template LaTeX files.

Please write a clear statement describing the problem you want to solve and the
motivation behind solving it. Find at least one recent (in the past 5 years)
published research paper that tackles this problem. You should approach your
project from a different angle than previous published works. Be creative!

\section{Data}
Please list and credit some datasets that you want to use for your project.
Explain how the dataset is related to the problem. \textbf{Make sure you can
download the datasets your report; otherwise you will have a hard time using
them!} Report some basic details regading your datasets, such as the size of
the dataset, the format (is it structured as phrase-image pairs? is it a corpus
of text?).

\section{Method}
Please list two or more directions on how you want to approach the problem. You
should mention machine learning techniques you plan on using such as SVM, PCA,
LDA, Neural Networks, feature design, etc... and you should explain the
benefits of your chosen techniques on the particular problem you are solving.

\section*{References}
Use unnumbered first-level heading for the references. Any choice of citation
style is acceptable as long as you are consistent. It is permissible to reduce
the font size to \verb+small+ (9 point) when listing the references.

\medskip

\small

[1] Alexander, J.A.\ \& Mozer, M.C.\ (1995) Template-based algorithms for
connectionist rule extraction. In G.\ Tesauro, D.S.\ Touretzky and T.K.\ Leen
(eds.), {\it Advances in Neural Information Processing Systems 7},
pp.\ 609--616. Cambridge, MA: MIT Press.

[2] Bower, J.M.\ \& Beeman, D.\ (1995) {\it The Book of GENESIS: Exploring
  Realistic Neural Models with the GEneral NEural SImulation System.}  New York:
TELOS/Springer--Verlag.

[3] Hasselmo, M.E., Schnell, E.\ \& Barkai, E.\ (1995) Dynamics of learning and
recall at excitatory recurrent synapses and cholinergic modulation in rat
hippocampal region CA3. {\it Journal of Neuroscience} {\bf 15}(7):5249-5262.

\end{document}
